\subsection{\texorpdfstring{空间 \(\mathcal{L}_{\mathfrak{S}}(\mathbf{E};\mathbf{F})\)}{空间 \textbackslash mathcal\{L\}\_\{\textbackslash mathfrak\{S\}\}(\textbackslash mathbf\{E\};\textbackslash mathbf\{F\})}}

设 \(\mathbf{F}\) 是一个拓扑向量空间，\(\mathbf{E}\) 是一个任意集合，\(\mathfrak{S}\) 是 \(\mathbf{E}\) 的子集的一个族。考虑向量空间 \(\mathrm{F}^{\mathrm{E}}\)，带上 \(\mathfrak{S}\) -收敛的一致结构（GT, X, § 1, No.~2）。我们知道这一结构与 \(\mathrm{F}^{\mathrm{E}}\) 的交换群结构相容（GT, X, § 1, No.~4, 推论 2）。如此得到的拓扑称为 \(\mathfrak{S}\) -拓扑。若 \(\mathbf{X}\) 是 \(\mathrm{F}^{\mathrm{E}}\) 的一个子集，或更一般地，是带有映射 \(j: \mathrm{X} \to \mathrm{F}^{\mathrm{E}}\) 的一个集合，则在 \(j\) 之下，\(\mathfrak{S}\) -拓扑（于 \(\mathrm{F}^{\mathrm{E}}\) 上）的原像称为 \(\mathfrak{S}\) -拓扑（在 \(\mathbf{X}\) 上）。

注. --- 1) \(\mathfrak{S}\) -拓扑与 \(\mathfrak{S}'\) -拓扑相同，其中 \(\mathfrak{S}'\) 表示由 \(\mathfrak{S}\) 生成的有界结构（III, 第 1 页）。

\begin{enumerate}
\def\labelenumi{\arabic{enumi})}
\setcounter{enumi}{1}
\tightlist
\item
  设 \(\mathbf{M} \in \mathfrak{S}\)，并设 \(\mathbf{V}\) 是 \(\mathbf{F}\) 中 0 的一个邻域；用 \(\mathrm{T}(\mathbf{M}, \mathbf{V})\) 表示所有这样的 \(f \in \mathrm{F}^{\mathrm{E}}\) 所成的集：\(f(x) \in \mathbf{V}\) 对每个 \(x \in \mathbf{M}\) 成立。若 \(\mathfrak{S}\) 在有限并下稳定，则诸集 \(\mathrm{T}(\mathbf{M}, \mathbf{V})\) 构成 \(\mathfrak{S}\) -拓扑（\(\mathrm{F}^{\mathrm{E}}\) 的）的 0 的一个基本邻域系。
\end{enumerate}

命题 1. --- 设 E 是一个集合，S 是 E 的子集的一个族，F 是一个拓扑向量空间，H 是 \(F^{E}\) 的一个向量子空间。为使 S-拓扑与 H 的向量空间结构相容，必要且充分的是 \(u(\mathbf{M})\) 对每个 \(u \in H\) 和每个 \(M \in S\) 都在 F 中有界。此外，若 F 是局部凸的，则 H 上的 S-拓扑是局部凸的。

根据上面的注 1) 和 2)，我们看到 \(\mathfrak{S}\) -拓扑与 H 的向量空间结构相容的一个必要且充分条件是：诸集 \(\mathrm{H} \cap \mathrm{T}(\mathbf{M}, \mathrm{V})\) 在 H 中是吸收的（I, 第 7 页, 命题 4）；但这蕴含着：对每个 \(u \in \mathrm{H}\)、每个子集 \(\mathbf{M} \in \mathfrak{S}\) 以及 F 中 0 的每个均衡邻域 V，存在 \(\lambda \neq 0\) 使得 \(u(M) \subset \lambda V\)；也就是说（III, 第 2 页），\(u(\mathbf{M})\) 在 F 中有界。最后，命题的最后一个断言由如下事实得出：若 V 是凸的，则 \(\mathrm{T}(\mathbf{M}, \mathrm{V})\) 也是凸的。

推论. --- 设 E 与 F 是两个局部凸空间，S 是 E 的有界子集的一个族，\(\mathcal{L}(E; F)\) 是从 E 到 F 中的连续线性映射所成的向量空间。则 S-拓扑与 \(\mathcal{L}(E; F)\) 的向量空间结构相容且是局部凸的。

只需指出：若 \(u\) 是从 \(\mathbf{E}\) 到 \(\mathbf{F}\) 中的一个连续线性映射，且 \(\mathbf{M}\) 是 \(\mathbf{E}\) 的一个有界子集，则 \(u(\mathbf{M})\) 在 \(\mathbf{F}\) 中有界（III, 第 4 页, 推论 1）。

给定两个局部凸向量空间 E 和 F 以及 E 的有界子集的一个族 S，用 \(\mathcal{L}_{\mathfrak{S}}(E;F)\) 表示给 \(\mathcal{L}(E;F)\) 赋予 S-拓扑所得到的局部凸空间。

例. --- 1) 若 \(\mathfrak{S}\) 是 E 的所有有限子集所成的集，则 \(\mathfrak{S}\) -拓扑是简单收敛拓扑，且空间 \(\mathcal{L}_{\mathfrak{S}}(E; F)\) 也记作 \(\mathcal{L}_s(E; F)\)。\(\mathcal{L}_s(E; F)\) 的一个有界子集称为 \(\mathcal{L}(E; F)\) 的一个简单有界子集。

\begin{enumerate}
\def\labelenumi{\arabic{enumi})}
\setcounter{enumi}{1}
\item
  若 \(\mathfrak{S}\) 是紧（相应地，预紧、紧凸）子集所成的集，则 \(\mathfrak{S}\) -拓扑称为紧（相应地，预紧、紧凸）收敛拓扑，且空间 \(\mathcal{L}_{\mathfrak{S}}(\mathrm{E};\mathrm{F})\) 也记作 \(\mathcal{L}_c(\mathrm{E};\mathrm{F})\)（相应地，\(\mathcal{L}_{pc}(\mathrm{E};\mathrm{F}),\mathcal{L}_{cc}(\mathrm{E};\mathrm{F})\)）。（参见 IV, 第 48 页, 习题 7。）
\item
  若 \(\mathfrak{S}\) 是 E 的所有有界子集所成的集，我们说 \(\mathfrak{S}\) -拓扑是有界收敛拓扑，且空间 \(\mathcal{L}_{\mathfrak{S}}(\mathrm{E};\mathrm{F})\) 记作 \(\mathcal{L}_b(\mathrm{E};\mathrm{F})\)。
\item
  当 \(\mathrm{F} = \mathbf{K}\) 时，空间 \(\mathcal{L}(\mathrm{E};\mathrm{F})\) 就是 \(\mathrm{E}'\)，也就是 \(\mathrm{E}\) 的对偶。我们用 \(\mathrm{E}_{\Xi}^{\prime}, \mathrm{E}_{s}^{\prime}\) 等表示空间 \(\mathcal{L}_{\Xi}(\mathrm{E};\mathbf{K}), \mathcal{L}_s(\mathrm{E};\mathbf{K})\) 等。空间 \(\mathrm{E}_s^\prime\)（相应地，\(\mathrm{E}_b^\prime\)）称为 \(\mathrm{E}\) 的弱对偶（相应地，强对偶）。\(\mathrm{E}_s^\prime\)（相应地，\(\mathrm{E}_b^\prime\)）的一个有界子集称为弱有界（相应地，强有界）的。我们注意到 \(\mathrm{E}'\) 上的弱拓扑不是别的，正是 \(\sigma (\mathrm{E}',\mathrm{E})\)（II, 第 42 页）。
\end{enumerate}

当 \(\mathrm{E} = \mathrm{F}\) 时，我们用 \(\mathcal{L}(\mathrm{E}), \mathcal{L}_{\mathfrak{S}}(\mathrm{E})\) 等表示空间 \(\mathcal{L}(\mathrm{E};\mathrm{F}), \mathcal{L}_{\mathfrak{S}}(\mathrm{E};\mathrm{F})\) 等。设 \(p\) 是 \(\mathrm{F}\) 上的一个连续半范数，\(\mathbf{M}\) 是 \(\mathrm{E}\) 的一个有界子集。设

\[
p _ {\mathrm{M}} (u) = \sup _ {x \in \mathrm{M}} p (u (x)).\tag{1}
\]

立即可见 \(p_{M}\) 是 \(\mathcal{L}(E;F)\) 上的一个半范数，并且若 \(\Gamma\) 是 F 上的一个基本半范数系，则半范数族 \(p_{M}\)（其中 p 取遍 \(\Gamma\)，M 取遍由 S 生成的有界结构的一个基）是 \(\mathcal{L}_{\mathfrak{S}}(E;F)\) 的一个基本半范数系。

特别地，若 E 和 F 是半赋范空间，且 p（相应地，q）表示 E（相应地，F）的半范数，则 \(\mathcal{L}(\mathrm{E};\mathrm{F})\) 上的有界收敛拓扑由半范数

\[
r (u) = \sup _ {p (x) \leqslant 1} q (u (x))\tag{2}
\]

（参见 GT, X, § 3, No.~2）。当我们把 \(\mathcal{L}_b(\mathrm{E};\mathrm{F})\) 看作半赋范空间时，除非另有明确说明，我们总是指半范数 (2)。若 F 是赋范空间，则半范数 (2) 是一个范数。

注. --- 3) 设 A 是 E 的单位球的一个稠密子集。鉴于 u 的连续性，我们还有

\[
r (u) = \sup _ {x \in \mathrm{A}} q (u (x)).\tag{3}
\]

例如

\[
r (u) = \sup _ {p (x) <   1} q (u (x)).\tag{4}
\]

由于 \(u(tx) = tu(x)\) 对 \(t \in \mathbf{R}\) 成立，我们还有，

\[
r (u) = \sup _ {p (x) = 1} q (u (x)) = \sup _ {p (x) \neq 0} \frac {q (u (x))}{p (x)}.\tag{5}
\]

只要 \(p \neq 0\)。

\begin{enumerate}
\def\labelenumi{\arabic{enumi})}
\setcounter{enumi}{3}
\tightlist
\item
  公式 (2) 表明映射 \(u \mapsto r(u)\) 在 \(\mathcal{L}_s(\mathrm{E};\mathrm{F})\) 上是下半连续的。
\end{enumerate}

命题 2. --- 设 E 和 F 是两个局部凸空间，并设 S 是 E 的有界子集所成的一个集。

\begin{enumerate}
\def\labelenumi{\arabic{enumi})}
\item
  \(\mathfrak{S}\) -拓扑（\(\mathcal{L}(\mathrm{E};\mathrm{F})\) 上的）与 \(\mathfrak{S}\) -拓扑相同，其中 \(\mathfrak{S}\) 表示 \(\mathrm{E}\) 上包含 \(\mathfrak{S}\) 的最小适配有界结构（III, 第 3 页）。
\item
  假设 \(\{0\}\) 在 \(\mathbf{F}\) 中不稠密，并设 \(\mathfrak{S}'\) 是 \(\mathbf{E}\) 的有界子集所成的另一个集。则 \(\mathfrak{S}'\) -拓扑粗于 \(\mathfrak{S}\) -拓扑，当且仅当 \(\mathfrak{S}' \subset \tilde{\mathfrak{S}}\)。
\end{enumerate}

设 \(u \in \mathcal{L}(\mathrm{E};\mathrm{F})\)，\(\mathbf{M} \in \mathfrak{S}\)，并设 \(p\) 是 \(\mathrm{F}\) 上的一个连续半范数。由于 \(p \circ u\) 是 \(\mathrm{E}\) 上的连续半范数，这就等于说 \(p \circ u\) 在 \(\mathbf{M}\) 上、或在闭凸均衡包 \(\tilde{\mathbf{M}}\)（\(\mathbf{M}\) 的）上以 1 为上界；换言之，我们有 \(p_{\mathbf{M}} = p_{\tilde{\mathbf{M}}}\)。此外，显然我们有 \(p_{\lambda \mathbf{M}} = \lambda p_{\mathbf{M}}\)（对所有 \(\lambda > 0\)），且有 \(p_{\mathbf{M} \cup \mathbf{M}'} = \sup(p_{\mathbf{M}}, p_{\mathbf{M}'})\)，由此即得第一个断言，因为 \(\tilde{\mathfrak{S}}\) 以 \(\mathfrak{S}\) 中集合的有限并的闭凸均衡包的位似所成之集为一个基。

我们现在证明第二个断言：首先，若 \(\mathbf{F}\) 是基域，则由定义可得，\(\tilde{\mathfrak{S}}\) -拓扑（在 \(\mathrm{E}' = \mathcal{L}(\mathrm{E};\mathrm{F})\) 上）以 \(\tilde{\mathfrak{S}}\) 中诸集的极集所成之集为 0 的一个基本邻域系。设 \(\mathbf{A}\) 是 \(\mathbf{E}\) 的一个有界子集，其极集 \(\mathbf{A}^{\circ}\) 是 \(\tilde{\mathfrak{S}}\) -拓扑下 0 的一个邻域；则存在闭凸均衡集 \(\mathbf{B} \in \tilde{\mathfrak{S}}\) 使得 \(\mathbf{A}^{\circ} \supset \mathbf{B}^{\circ}\)，于是 \(\mathbf{A} \subset \mathbf{B}^{\circ \circ}\)；但由 II, 第 45 页 的推论 3，有 \(\mathbf{B}^{\circ \circ} = \mathbf{B}\)，因此 \(\mathbf{A} \subset \mathbf{B}\) 且 \(\mathbf{A} \in \tilde{\mathfrak{S}}\)。于是，若 \(\mathfrak{S}'\) 是 \(\mathbf{E}\) 的有界子集所成的一个集，则 \(\mathfrak{S}'\) -拓扑粗于 \(\mathfrak{S}\) -拓扑（\(\mathrm{E}'\) 上的），当且仅当 \(\mathfrak{S}' \subset \tilde{\mathfrak{S}}\)。一般情形立即可得，因为若 \(y \in F\) 不在 0 的闭包中，则可以验证：使 \(f \in E'\) 对应于映射 \(x \mapsto f(x)y\) 的映射是局部凸空间 \(\mathrm{E}'_{\mathfrak{S}}\) 到它在 \(\mathcal{L}_{\mathfrak{S}}(\mathrm{E};\mathrm{F})\) 中的像上的一个同构。

